\documentclass[11pt]{article}

\usepackage[margin=1in]{geometry}
\usepackage{amsmath,amssymb}
\usepackage{booktabs}
\usepackage{array}
\usepackage{enumitem}
\usepackage{hyperref}

\setlist[enumerate]{itemsep=0.4em}
\setlist[itemize]{itemsep=0.25em}
\renewcommand{\arraystretch}{1.18}

\title{STAT 514 Practice Final Solutions}
\author{}
\date{}

\begin{document}
\maketitle

\noindent Source: \texttt{midterms/midterm3/practice.pdf}. These notes emphasize reusable formulas and exam-ready calculations.

\section*{Problem 1}

\begin{enumerate}[label=(\alph*)]
\item A \(2^{8-3}_{IV}\) design has
\[
2^{8-3}=32
\]
runs and \(8\) two-level factors.

\item In a resolution IV design, a two-factor interaction is not aliased with a main effect, but it may be aliased with another two-factor interaction.

\item A \(2^{9-5}_{III}\) design has \(2^{9-5}=16\) runs, so it can estimate
\[
16-1=15
\]
factorial effects besides the intercept. Since \(p=5\), each effect is aliased with
\[
2^p-1=2^5-1=31
\]
other factorial effects.

\item The fundamental principles for factorial effects are sparsity, hierarchy, and heredity.

\item A CCD using a \(2^{5-1}_V\) factorial portion and \(3\) center points has
\[
n_F+2k+n_C=2^{5-1}+2(5)+3=16+10+3=29
\]
runs.

\item A \(3^2\) full factorial design has \(3^2=9\) runs and \(2\) factors.

\item A Latin hypercube design with \(n\) runs in \(k\) factors has \(1\) point in each of the \(n\) equally sized intervals of every one-dimensional marginal, so each level is sampled \(1\) time per factor.

\item A maximin LHD maximizes the minimum pairwise distance. A minimax design minimizes the maximum distance from any candidate input to its nearest design point.
\end{enumerate}

\section*{Problem 2}

General rule:
\[
df_E=N-p,
\]
where \(p\) includes the intercept and all fitted model parameters.

\begin{enumerate}[label=(\alph*)]
\item CRD with three treatments and \(7,8,9\) replicates:
\[
N=24,\qquad df_E=N-a=24-3=21.
\]

\item Paired comparison with \(15\) pairs:
\[
df_E=n_{\text{pairs}}-1=15-1=14.
\]

\item Full \(2^4\) factorial with \(3\) replicates and all effects:
\[
N=2^4(3)=48,\qquad p=2^4=16,\qquad df_E=48-16=32.
\]
Equivalently, \(df_E=2^k(n-1)=2^4(2)=32\).

\item Unreplicated \(2^{8-2}\) design:
\[
N=2^6=64.
\]
The model has
\[
p=1+8+\binom{8}{2}=1+8+28=37.
\]
Thus
\[
df_E=64-37=27.
\]

\item CCD using \(2^{3-1}\) factorial runs and \(5\) center runs:
\[
N=2^{3-1}+2(3)+5=4+6+5=15.
\]
A standard second-order model for \(k=3\) has
\[
p=1+k+k+\binom{k}{2}=1+3+3+3=10.
\]
Thus
\[
df_E=15-10=5.
\]

\item Full \(3^2\) factorial with two replicates:
\[
N=3^2(2)=18,\qquad p=9,
\]
so
\[
df_E=18-9=9.
\]
\end{enumerate}

\section*{Problem 3}

\begin{enumerate}[label=(\alph*)]
\item Use a \(16\)-run \(2^{7-3}\) fractional factorial design, preferably maximum resolution IV. This is appropriate for screening seven two-level factors under a 16-run limit. There is no stated blocking factor. Sample size:
\[
N=16.
\]

\item Use a central composite design for the three quantitative factors because a second-order model is needed for curvature, interactions, prediction, and optimization. No stated blocking factor. General sample size:
\[
N=2^3+2(3)+n_C=14+n_C.
\]
With five center runs:
\[
N=19.
\]

\item Use a space-filling design, preferably a maximin Latin hypercube design, for \(12\) simulator inputs. The simulator is deterministic, so replication is wasteful. No blocking factor. Since the scientists want 100 uniformly spaced levels per parameter:
\[
N=100.
\]

\item Use a completely randomized design. The aquaria are essentially homogeneous and there is no obvious nuisance factor. No blocking factor. Balanced allocation gives:
\[
N=40,\qquad 10\ \text{aquaria per diet}.
\]

\item Use a randomized complete block design with production day as the block. Each of the three treatment levels is run within each of five days:
\[
N=3(5)=15.
\]
\end{enumerate}

\section*{Problem 4}

\subsection*{(a)}
Let \(i=1,2,3\) index cutting speed, \(j=1,2,3\) index coolant type, and \(k=1,\ldots,4\) index replicates. The two-factor ANOVA model with interaction is
\[
y_{ijk}=\mu+\alpha_i+\beta_j+(\alpha\beta)_{ij}+\varepsilon_{ijk},
\qquad
\varepsilon_{ijk}\stackrel{iid}{\sim}N(0,\sigma^2).
\]
The constraints are
\[
\sum_{i=1}^3\alpha_i=0,\qquad
\sum_{j=1}^3\beta_j=0,
\]
\[
\sum_{i=1}^3(\alpha\beta)_{ij}=0\quad\text{for each }j,\qquad
\sum_{j=1}^3(\alpha\beta)_{ij}=0\quad\text{for each }i.
\]

\subsection*{(b)}
There are \(N=3\cdot3\cdot4=36\) observations. Degrees of freedom are
\[
df_A=2,\quad df_B=2,\quad df_{AB}=4,\quad df_E=3\cdot3(4-1)=27,\quad df_T=35.
\]
Given \(R^2=0.94\) and \(SSE=18\),
\[
0.94=1-\frac{18}{SST}\quad\Rightarrow\quad SST=\frac{18}{0.06}=300.
\]
Thus
\[
MSE=\frac{18}{27}=\frac{2}{3}.
\]
Cutting speed has \(F=90\):
\[
MS_A=90\left(\frac{2}{3}\right)=60,\qquad SS_A=2(60)=120.
\]
Coolant type has \(SS_B=90\):
\[
MS_B=45,\qquad F_B=\frac{45}{2/3}=67.5.
\]
Interaction:
\[
SS_{AB}=300-120-90-18=72,\quad
MS_{AB}=\frac{72}{4}=18,\quad
F_{AB}=\frac{18}{2/3}=27.
\]

\[
\begin{array}{lrrrr}
\toprule
\text{Source} & df & SS & MS & F\\
\midrule
\text{Cutting speed} & 2 & 120 & 60 & 90\\
\text{Coolant type} & 2 & 90 & 45 & 67.5\\
\text{Interaction} & 4 & 72 & 18 & 27\\
\text{Error} & 27 & 18 & 2/3 & \\
\text{Total} & 35 & 300 & & \\
\bottomrule
\end{array}
\]

\section*{Problem 5}

For a one-replicate \(2^k\) factorial, the effect of term \(T\) is
\[
\widehat{\text{effect}}_T=\bar y_{T+}-\bar y_{T-}
=\frac{1}{2^{k-1}}\sum_{\text{runs}}s_Ty,
\]
where \(s_T\) is the product of coded signs for the term.

\begin{enumerate}[label=(\alph*)]
\item Main effect of \(A\):
\[
\bar y_{A+}=\frac{72+85+79+80}{4}=79,\qquad
\bar y_{A-}=\frac{98+87+99+87}{4}=92.75.
\]
\[
\hat A=79-92.75=-13.75.
\]

\item For \(AB\), \(s_{AB}=x_Ax_B\):
\[
\bar y_{AB+}=\frac{98+85+99+80}{4}=90.5,\qquad
\bar y_{AB-}=\frac{72+87+79+87}{4}=81.25.
\]
\[
\widehat{AB}=90.5-81.25=9.25.
\]

\item Main effect plot for \(A\): plot \((-1,92.75)\) and \((+1,79.00)\). The line decreases, so high \(A\) lowers roughness.

For the \(AB\) interaction plot:
\[
\begin{array}{crr}
\toprule
B & \bar y(A=-1) & \bar y(A=+1)\\
\midrule
-1 & (98+99)/2=98.5 & (72+79)/2=75.5\\
+1 & (87+87)/2=87.0 & (85+80)/2=82.5\\
\bottomrule
\end{array}
\]
The lines are not parallel; the effect of \(A\) is stronger when \(B=-1\).

\item Reduced model:
\[
y=\beta_0+\beta_Ax_A+\beta_{AB}x_Ax_B+\varepsilon.
\]
For an orthogonal two-level factorial:
\[
\hat\beta_0=\bar y,\qquad \hat\beta_T=\frac{\widehat{\text{effect}}_T}{2}.
\]
Thus
\[
\hat\beta_0=\frac{98+72+87+85+99+79+87+80}{8}=85.875,
\]
\[
\hat\beta_A=-13.75/2=-6.875,\qquad
\hat\beta_{AB}=9.25/2=4.625.
\]
Fitted model:
\[
\hat y=85.875-6.875x_A+4.625x_Ax_B.
\]

\item Predictions:
\[
\begin{array}{rrrr}
\toprule
A & B & x_Ax_B & \hat y\\
\midrule
-1 & -1 & +1 & 97.375\\
+1 & -1 & -1 & 74.375\\
-1 & +1 & -1 & 88.125\\
+1 & +1 & +1 & 83.625\\
\bottomrule
\end{array}
\]
The minimum predicted roughness is at \(A=+1,B=-1\), with \(\hat y=74.375\).
\end{enumerate}

\section*{Problem 6}

\begin{enumerate}[label=(\alph*)]
\item The column \(D\) equals \(ABC\), so
\[
D=ABC,\qquad I=ABCD.
\]

\item The column \(E\) equals \(-AC\), so
\[
E=-AC,\qquad I=-ACE.
\]

\item The independent words are \(I=ABCD\) and \(I=-ACE\). Multiplying them gives
\[
(ABCD)(-ACE)=-BDE.
\]
Therefore the complete defining relation is
\[
I=ABCD=-ACE=-BDE.
\]

\item The shortest nonidentity words are \(ACE\) and \(BDE\), both of length 3. The design has resolution III.

\item Reverse the signs of column \(B\). If the reversed column is still called \(B\), then
\[
D=-ABC,\qquad E=-AC.
\]
Thus
\[
I=-ABCD,\qquad I=-ACE.
\]
Multiplying gives \(BDE\), so the complete defining relation is
\[
I=-ABCD=-ACE=BDE.
\]
\end{enumerate}

\section*{Problem 7}

\begin{enumerate}[label=(\alph*)]
\item
\[
\bar y_{A+}=\frac{32+45}{2}=38.5,\qquad
\bar y_{A-}=\frac{54+47}{2}=50.5.
\]
\[
\hat A=38.5-50.5=-12.
\]

\item
\[
\bar y_{AB+}=\frac{54+45}{2}=49.5,\qquad
\bar y_{AB-}=\frac{47+32}{2}=39.5.
\]
\[
\widehat{AB}=49.5-39.5=10.
\]

\item Estimate pure error from the center points \(40,38,42\):
\[
\bar y_C=40,
\]
\[
\hat\sigma^2=MSE_{\text{pure error}}
=\frac{(40-40)^2+(38-40)^2+(42-40)^2}{3-1}
=4.
\]
Degrees of freedom:
\[
df_E=3-1=2.
\]

\item For a one-replicate \(2^2\) factorial,
\[
SE(\widehat{\text{effect}})=\sqrt{MSE}.
\]
Thus
\[
SE(\widehat{AB})=\sqrt{4}=2,\qquad
t_0=\frac{10}{2}=5.
\]
With \(df=2\), \(t_{0.025,2}=4.303\). Since \(|5|>4.303\), the \(AB\) interaction is significant at the 5\% level.

\item Curvature compares center and factorial means:
\[
\bar y_F=\frac{54+47+32+45}{4}=44.5,\qquad
\bar y_C=40.
\]
\[
\hat C=\bar y_C-\bar y_F=-4.5.
\]
Standard error:
\[
SE(\hat C)=\sqrt{MSE\left(\frac{1}{n_C}+\frac{1}{n_F}\right)}
=\sqrt{4\left(\frac{1}{3}+\frac{1}{4}\right)}
=\sqrt{\frac{7}{3}}=1.528.
\]
Test statistic:
\[
t_0=\frac{-4.5}{1.528}=-2.95.
\]
Since \(|-2.95|<t_{0.025,2}=4.303\), curvature is not significant at the 5\% level.

Equivalent:
\[
SS_{\text{curv}}=\frac{n_Fn_C}{n_F+n_C}(\bar y_F-\bar y_C)^2
=\frac{4\cdot3}{7}(4.5)^2=34.714,
\]
\[
F_{\text{curv}}=\frac{34.714}{4}=8.679=t_0^2.
\]
\end{enumerate}

\section*{Problem 8}

The fitted model is
\[
y=8+4x_1-2x_2-2x_1^2-x_2^2+\varepsilon.
\]

\begin{enumerate}[label=(\alph*)]
\item A circumscribed CCD with \(k=2\), \(\alpha=\sqrt{2}\), and three center points has:
\[
\text{factorial points: }(-1,-1),(1,-1),(-1,1),(1,1),
\]
\[
\text{axial points: }(\sqrt2,0),(-\sqrt2,0),(0,\sqrt2),(0,-\sqrt2),
\]
\[
\text{center points: }(0,0),(0,0),(0,0).
\]
Total:
\[
N=4+4+3=11.
\]

\item In
\[
y=\beta_0+x'b+x'Bx+\varepsilon,\qquad x=(x_1,x_2)',
\]
we have
\[
\beta_0=8,\qquad
b=\begin{pmatrix}4\\-2\end{pmatrix},
\]
and, because there is no interaction term,
\[
B=\begin{pmatrix}-2&0\\0&-1\end{pmatrix}.
\]

\item The stationary point solves
\[
b+2Bx_s=0,\qquad x_s=-\frac12B^{-1}b.
\]
Since
\[
B^{-1}=\begin{pmatrix}-1/2&0\\0&-1\end{pmatrix},
\]
\[
B^{-1}b=\begin{pmatrix}-2\\2\end{pmatrix},
\qquad
x_s=-\frac12\begin{pmatrix}-2\\2\end{pmatrix}
=\begin{pmatrix}1\\-1\end{pmatrix}.
\]
The stationary point is \((x_1,x_2)=(1,-1)\).

\item The eigenvalues of the diagonal matrix \(B\) are
\[
\lambda_1=-2,\qquad \lambda_2=-1.
\]
Both are negative, so \(B\) is negative definite and the stationary point is a maximum. The predicted maximum response is
\[
\hat y(1,-1)=8+4(1)-2(-1)-2(1)^2-(-1)^2=11.
\]
\end{enumerate}

\end{document}
